\subsection{Existing solutions to the index number problem}
\label{sec:existing_index_number}
A few proposed methods exist that compute decompositions similar to the (N)OBD, and yet do not have the index number problem.
We describe the approach of these options here and discuss how they do not resolve the issue of arbitrariness discussed in this paper.

\textbf{Shapley Values.} Shapley values \citep{Shapley1953} are widely used in the machine learning literature to understand a models' predictions.
A straightforward use of them would be to estimate ``explained'' or ``unexplained'' contributions to $\Delta_Y$.
\textcolor{red}{This follows from the standard Shapley construction, which treats each reference-based decomposition as a ``player'' and assigns contributions by averaging marginal effects over all orderings; with two ``players'' (the two reference directions $H$ and $K$), there are two equally weighted permutations, $(H,K)$ and $(K,H)$, so the Shapley value reduces to the simple average}, for the explained component, $E = \tfrac{1}{2}E^{(H)} + \tfrac{1}{2}E^{(K)}$.
However, one might see our paper as questioning the meaningfulness of $E^{(H)}$ and $E^{(K)}$. Then the Shapley proposal averages two questionable quantities!

\textbf{Pooled approaches.}
Another approach to eliminate the index number problem is to define some ``pooled'' group and then measure deviations from this group in a three-term decomposition. For example, in the case of the (linear) OBD, one defines a pooled $\beta^*$ and then computes:
\begin{equation}
  \mathbb{E}_H[Y] - \mathbb{E}_K[Y] = (\mu_H - \mu_K)^T \beta^* + \mu_H^T (\beta_H - \beta^*) + \mu_K^T(\beta^* - \beta_K).
\end{equation}
The first term is interpreted as differences due differences in covariates, the second differences due to differences in $H$'s $Y|X$ from the baseline, and the third differences due to differences in $K$'s $Y|X$ from the baseline.

The question that remains, then is how to define $\beta^*$.
A typical choice is to use some sort of ``average'' of the two groups.
One approach from \citet{oaxacaRansom1998}, defines some $w \in [0,1]$ and sets $\beta^* = (1-w)\beta_K + w\beta_H$.
A second approach, from \citet{neumark1988}, is to create a pooled dataset consisting of all samples from $H$ and $K$, and estimate $\beta^*$ from this dataset.
However, these approaches suffer from issues of arbitrariness as well.
First, how does one pick $w$ in the Oxaca-Ransom approach?
Second, why is defining $\beta^*$ one way more reasonable than the other?
If we assume both are equally reasonable, then, given our results for the OBD, we feel that one should be highly suspicious of substantive conclusions being robust to the choice of $\beta^*$.
We leave experimentation along these lines for future work.

\textbf{Intermediate covariate distributions.}
\citet{dinardoFortinLemieux1996} propose a variant of the OBD with a more fine-grained decomposition of the explained component, and \citet{quintero2025} propose a NOBD with a similarly fine-grained explained decomposition (\citet{quintero2025} also decomposes the unexplained component).
Both cases define a number of intermediate distributions over covariates $P_1, \dots, P_I$, where $P_1 = K_X$ and $P_I = H_X$. E.g., in the case of \citet{quintero2025}, $P_2$ is equal to $K$, except that the marginal of the first covariate is equal to the marginal of $H_X$.
In either case, for each $i = 2, \dots, I$, intermediate explained components are defined as one of:
\begin{align}
  \mathbb{E}_{X \sim P_i}[ \mathbb{E}_{Y \sim H}[Y|X] ] - \mathbb{E}_{X \sim P_{i-1}} [ \mathbb{E}_{Y \sim H} [Y|X]] \\
  \mathbb{E}_{X \sim P_i}[ \mathbb{E}_{Y \sim K}[Y|X] ] - \mathbb{E}_{X \sim P_{i-1}} [ \mathbb{E}_{Y \sim K} [Y|X]].
\end{align}  
There are two issues with such approaches.
First, \citet{dinardoFortinLemieux1996} does not change how the the unexplained component is computed over the OBD, so the same exact index number problems with the unexplained component detailed here apply to the decomposition of \citet{dinardoFortinLemieux1996}.
Second, one now has to justify not just a choice of $H$ or $K$ as a reference group, but instead a whole ordering of the $P_i$.
This problem is explicitly called out by \citet{quintero2025}, who choose to fix an arbitrary order.
Given our results here, we strongly suspect that conclusions will not be robust to the ordering of the $P_i$, and we leave such investigation for future work.
